\SourceLeaf{58}
\fancyhead[L]{\small ANZEIGEN EIGNER SCHRIFTEN}
\pdfbookmark[1]{Theoria motus corporum coelestium — Anzeige 1809}{anzeige-theoria1809}
\begin{center}\rule{.75\textwidth}{.3pt}\\
Göttingische gelehrte Anzeigen. 1809. Junius 17.\\
\rule{.75\textwidth}{.3pt}\end{center}
\textit{Theoria motus corporum coelestium in sectionibus conicis Solem ambientium.} Auctore \textsc{Carolo Frid. Gauss.} Hamburgi 1809. Sumtibus \textit{Frid. Perthes} et \textit{J. H. Besser.} XII S. Vorrede, 228 S. Text und 20 S. Tabellen, nebst einer Kupfertafel. gr. Quart.

Mit demjenigen Theile der theorischen Astronomie, welcher die parabolische und elliptische Bewegung der Himmelskörper zum Gegenstande hat, haben sich bekanntlich viele Schriftsteller, und unter ihnen sogar Geometer vom ersten Range in eigenen Werken, beschäftigt: nach solchen Vorgängern durfte nur eine grosse Veranlassung eine neue Bearbeitung dieses Feldes motiviren. Dem eigentlichen Astronomen ist diese Veranlassung bekannt genug: nur solchen Freunden der Himmelskunde, die aus derselben kein Hauptgeschäft machen, und vielleicht, in den letzten verhängnissvollen Jahren, über die Angelegenheiten der Erde die Angelegenheiten des Himmels aus dem Gesicht verloren haben möchten, wollen wir mit Wenigem die Umstände in Erinnerung bringen, welche zunächst gegenwärtiges Werk veranlasst haben. Die Aufgabe, aus den nur eine mässig lange Zeit hindurch von der Erde aus beobachteten Bewegungen eines Himmelskörpers, von dem man nichts weiter weiss, als dass er in einem Kegelschnitte nach den \textsc{Kepler}schen Gesetzen sich um die Sonne bewegt, dessen Bahn mit hinreichender Genauigkeit zu bestimmen, war bisher eigentlich noch nie auf eine
\SourceLeaf{59}ernstliche Art bearbeitet. Allerdings sind einige auf dieses Problem Bezug habende Untersuchungen vorhanden: allein ohne den Scharfsinn und die analytische Kunst zu verkennen, wovon einige derselben Spuren zeigen, muss man den vorgeschlagenen Methoden doch alles absprechen. was zur wirklichen Brauchbarkeit erfordert wird, indem sie entweder statt möglich grösster Schärfe nur höchstens eine rohe Annäherung geben, oder statt eines geschmeidigen, für die wirkliche Anwendung geformten, Calculs nur einen verworrenen Haufen von unentwickelten und selbst den unverdrossensten Rechner zurückschreckenden Formeln aufstellen, oder endlich statt auf Beobachtungen, wie sie der heutige Zustand der practischen Astronomie erlaubt, anwendbar zu sein, selbst schon durch weit kleinere Fehler, als diejenigen, welche bei den Beobachtungen unvermeidlich sind, ganz unbrauchbar werden. Der Grund dieser Vernachlässigung eines Problems, welches unstreitig schon an sich von einem hohen Interesse ist, scheint zum Theil in dem Umstande zu liegen, dass diejenigen Geometer, welche sich mit jenem Problem beschäftigten, mit den Kräften und Bedürfnissen der Ausübung nicht vertraut genug waren, hauptsächlich aber wohl darin, dass die Geschichte der Astronomie noch keinen Fall aufgestellt hatte, wo das Bedürfniss einer angemessenen Auflösung der Aufgabe recht dringend, und ihr Nutzen recht fühlbar gewesen wäre. In der That, als \textsc{Kepler} nach Entdeckung seiner Gesetze die Bestimmung der Dimensionen der Bahnen der damals bekannten Planeten unternahm, stand ihm, ausser den schon sehr genau bekannten mittlern Bewegungen, ein Schatz von guten und vieljährigen \textsc{Tycho}nischen Beobachtungen zu Gebote, aus welchen er nur auswählen durfte, was er zur Anwendung seiner zwar schönen, aber doch speciellen, und, verhältnissmässig, kunstlosen Methoden jedesmal nöthig fand. Dieselben Hülfsquellen, oder vielmehr noch grössere, hatten \textsc{Kepler}'s Nachfolger, die bei dem Fortschreiten der Beobachtungskunst die schon ziemlich nahe bekannten Elemente der Planetenbahnen noch genauer zu bestimmen unternahmen.

Anders verhielt es sich bei den Cometen, in deren Bewegung man nach \textsc{Newton}'s Entdeckungen nur einen besondern Fall der allgemeinen \textsc{Kepler}'schen Gesetze erkannte. Diese Weltkörper sind gewöhnlich nur eine kurze Zeit sichtbar: die Erscheinung ihrer Bewegungen hängt von den zufälligen Stellungen ab, welche die Erde während jener Sichtbarkeit eingenommen hatte, und der Geometer, welcher die Bestimmung der Bahnen aus den Erscheinungen unternimmt,
\SourceLeaf{60}kann dazu unter den Beobachtungen nicht viel aussuchen, er muss sie nehmen, wie er sie vorfindet, wie sie der Zufall gegeben hat, nur höchst selten stehen ihm Beobachtungen zu Gebote, die unter solchen Umständen gemacht sind, wie er sie zur Anwendung von speciellen Methoden verlangen möchte. Das Problem also, aus einigen Beobachtungen eines Cometen dessen parabolische Bahn zu bestimmen, war eben so wichtig, als schwer. \textsc{Newton} selbst erkannte die Schwierigkeiten an, und wusste sie zu besiegen; mit welchem Eifer und mit welchem Erfolge man seit \textsc{Newton} bis zu unsrer Zeit sich mit dieser Aufgabe beschäftigt hat, ist bekannt genug. Allein zwischen dieser Aufgabe und der oben von uns erwähnten findet ein wesentlicher Unterschied Statt: in jener wird die Art des Kegelschnitts \textit{vorgeschrieben}, während sie in diesem aus der unendlichen Mannigfaltigkeit möglicher Kegelschnitte, von denen die Parabel nur Eine Art ist, ohne hypothetische Voraussetzungen, blos mit Hülfe der Beobachtungen selbst, ausgemittelt werden soll. Begreiflich erhält jene Aufgabe durch diese Einschränkung eine grosse Vereinfachung, welche man sich aber bei den meisten Cometen erlauben durfte, und erlauben musste, da der gewöhnlich ziemlich geringe Grad von Genauigkeit in den Beobachtungen und ihre kurze Dauer kaum jemals hinreichen. das Dasein einer Abweichung von der Parabel zu beweisen, und ihre Grösse zu bestimmen. Freilich hat man doch bei einigen Cometen diese Bestimmung wirklich versucht, allein, und dies ist wesentlich, immer erst, nachdem man schon eine parabolische Bahn berechnet hatte, die dann als Annäherung zu der zu bestimmenden Ellipse oder Hyperbel diente. So blieb also auch hier unser allgemeineres Problem gewissermassen entbehrlich.

Auch die Entdeckung eines neuen perennirenden Weltkörpers im Jahre 1781 machte das Bedürfniss einer Auflösung dieses Problems noch nicht fühlbar. Nachdem man die Unzulänglichkeit einer parabolischen Bahn eingesehen hatte, versuchte man einen Kreis, der bei der zufälliger Weise ziemlich kleinen Excentricität der Bahn des Planeten seine Bewegung während einiger Jahre erträglich genau darstellte; bei der so sehr langsamen Bewegung des Planeten, seiner geringen Entfernung von der Ecliptik und seinem noch ziemlich lebhaften Lichte war hier auch weiter kein \textit{periculum in mora}; ohne alle Mühe fand man ihn von einem Jahre zum andern wieder auf, und zur Bestimmung der Abweichung vom Kreise oder der wahren Ellipse durfte man also warten, bis man die Beobachtungen nach seiner Bequemlichkeit aussuchen konnte.

\SourceLeaf{61}Ganz anders aber verhielt es sich mit der im Jahr 1801 entdeckten Ceres. Dieser Weltkörper zeigt sich nur als Sternchen achter Gröse, ist nur mit Mühe und bei genauer Kenntniss seines jedesmaligen Platzes aus dem zahllosen Heere ganz ähnlicher Fixsterne herauszufinden: der Entdecker hatte ihn nur während des kurzen Zeitraums von 41 Tagen beobachtet, und als die Entdeckung in dem übrigen Europa bekannt wurde, war er bereits in den Sonnenstrahlen verloren, um erst nach einem Jahre in einer ganz verschiedenen Himmelsgegend wieder sichtbar zu werden. Jetzt galt es die möglich genaueste Vorhersagung des Orts, wo man ihn wieder zu suchen haben würde, und diese musste bloss auf die wenigen vorhandenen Beobachtungen und strengen Calcul, ohne unsichere Hypothesen, gegründet werden. Mehrere Astronomen versuchten die einfachste Hypothese, eine Kreisbahn, mit der sich die Beobachtungen freilich nur in eine unvollkommene Übereinstimmung bringen liessen, und deren Zulänglichkeit zur Wiederauffindung also wenigstens sehr precär blieb: in der That hat der Erfolg nachher bestätigt, dass diese Kreishypothese schon am Ende des Jahres 1801 um eilf Grade von dem wahren Orte des Planeten abwich, und diejenigen Astronomen, welchen das Glück zu Theil ward, denselben wieder aufzufinden, haben selbst erklärt, dass diese Wiederauffindung nach einer so fehlerhaften Hypothese unmöglich gewesen sein würde.

Dem Verfasser des vorliegenden Werks hatten sich im Sommer 1801 bei Gelegenheit einer ganz andern Beschäftigung einige Ideen dargeboten, die ihm zu einer Auflösung des erwähnten allgemeinen Problems führen zu können schienen. Zu einer andern Zeit würde er vielleicht diese Ideen, welche zunächst nur theoretischen Reiz für ihn hatten, nicht sogleich weiter verfolgt und ausgeführt haben: allein gerade in jenem Zeitpunkte, wo \textsc{Piazzis} Entdeckung die allgemeine Aufmerksamkeit gespannt hatte, und dessen Beobachtungen so eben ins Publicum gekommen waren, konnte er sich nicht enthalten, an diesen die practische Anwendbarkeit jener Ideen zu prüfen. Der Erfolg dieser Arbeit ist bekannt. Die bis dahin nicht geahnte Möglichkeit, aus einer kurzen Reihe von Beobachtungen eines Planeten eine schon sehr genäherte und zu seiner Wiederauffindung nach einem grössern Zeitraume überflüssig genaue Bestimmung seiner Bahn zu machen, war dadurch aufs schönste erwiesen, und die Brauchbarkeit der angewandten Methode bewährt: und wenn über die Allgemeinheit dieser Brauchbarkeit noch Zweifel hätten übrig bleiben können, so sind diese durch eben so
